\chapter{Poisson Brackets, Canonical Transformations, and Symmetry}

We continue the discussion of Poisson brackets, but now in the wider language of structure and flow. The point of departure is not a computational trick, but the claim that phase space itself carries a special invariant structure, and that the mathematics of classical mechanics is the mathematics of transformations that respect it. These notes follow Leonard Susskind's Lecture 9 closely, with transcript and frame support preserved where the boardwork matters; the lecture is due to Leonard Susskind, and the supporting curation used for these notes comes through LazyingArt LLC and Video2Book.

\section{Poisson Structure and the Idea of Flow}

Phase space is not merely a list of coordinates. It is not just a collection of \(q\)'s and \(p\)'s written side by side. It is a space with structure, and the relevant structure is the one left invariant by the transformations we are allowed to perform.

A familiar example of structure is metric structure. On a Riemannian space one has a notion of distance between nearby points, and that metric data is the geometric structure of the space. The Poisson structure of phase space is different. It is more abstract, and it is tailored not to distance but to the algebraic relations among phase-space coordinates and to the flows that move points through phase space.

This is the first important shift in viewpoint. Rotations and translations in ordinary coordinate space may already be thought of as flows: one imagines passing continuously through a family of slightly rotated axes, or through a family of slightly translated coordinate systems. The same idea can be enlarged. In phase space we may contemplate flows that mix coordinates and momenta. Those richer flows are precisely what Poisson brackets are built to describe.

One kind of flow is the actual time development of a system under a Hamiltonian. But it is not the only kind. A symmetry may also be understood as a flow, even when it is not the physical time evolution of the system. This is why the lecture insists on starting from the structure of phase space before returning to formulas.

\section{Abstract Rules for Poisson Brackets}

At this point the lecture deliberately backs up and abstracts. Instead of beginning with the derivative definition, we begin with the algebraic properties. Let \(A\), \(B\), and \(C\) be ordinary commuting functions on phase space. Then the Poisson bracket satisfies

\begin{equation}
\{A,B\}=-\{B,A\},
\end{equation}

\begin{equation}
\{\alpha A+\beta B,C\}=\alpha\{A,C\}+\beta\{B,C\},
\end{equation}

and the product rule
\begin{equation}
\{AB,C\}=A\{B,C\}+B\{A,C\}.
\end{equation}

The point about ordinary multiplication is worth keeping in view. The functions \(A\), \(B\), and \(C\) commute as numbers do, so the order of multiplication outside the bracket is immaterial in classical mechanics. What does matter is the order inside the bracket. Reversing the order of the bracket reverses the sign.

We also assume, as the lecture does, that the functions under discussion are smooth enough for differentiation and repeated differentiation to make sense. When we do need the concrete definition, we return to the familiar formula
\begin{equation}
\{A,B\}=\sum_i \left(\frac{\partial A}{\partial q_i}\frac{\partial B}{\partial p_i}-\frac{\partial A}{\partial p_i}\frac{\partial B}{\partial q_i}\right).
\end{equation}
But the narrative order matters: first the structural rules, then the concrete computations.

\section{Canonical Brackets and Bracketing as Differentiation}

The lecture next descends from the abstract rules to the concrete brackets of the canonical variables themselves. At this stage the board uses capital \(Q\)'s and \(P\)'s abstractly for canonical coordinates. Later, when transformed coordinates enter, we will reserve capitals for the transformed variables and return to lower-case \(q\)'s and \(p\)'s for the original ones.

The surviving board frame catches exactly the moment when the abstract product rule sits above the simplest concrete consequence.

\begin{figure}[h]
\centering
\includegraphics[width=0.56\textwidth]{lecture_09_figure_02.png}
\caption{Poisson bracket of \(P\) with itself}
\end{figure}

The upper line in the frame is cropped, but the context is the product rule already written above. The lower line isolates the special case
\begin{equation}
\{P,P\}=0.
\end{equation}
The same argument gives
\begin{equation}
\{Q,Q\}=0.
\end{equation}

For many degrees of freedom, the concrete canonical relations are
\begin{equation}
\{Q_i,Q_j\}=0,\qquad \{P_i,P_j\}=0,\qquad \{Q_i,P_j\}=\delta_{ij},
\end{equation}
and therefore
\begin{equation}
\{P_i,Q_j\}=-\delta_{ij}.
\end{equation}

From the derivative definition one also sees immediately that any two functions depending only on the \(Q\)'s have vanishing Poisson bracket, and likewise for two functions depending only on the \(P\)'s:
\begin{equation}
\{F(Q),G(Q)\}=0,\qquad \{F(P),G(P)\}=0.
\end{equation}
To get anything nonzero, one needs a \(Q\) on one side to meet a \(P\) on the other.

\subsection{Question \& Answer}

Question. Why does antisymmetry immediately force \(\{P,P\}=0\), but not yet determine \(\{Q,P\}\)?

Answer. For a self-bracket, antisymmetry gives
\[
\{P,P\}=-\{P,P\},
\]
so the bracket equals its own negative and must vanish. But for \(\{Q,P\}\), antisymmetry only says
\[
\{Q,P\}=-\{P,Q\}.
\]
That is a relation between two different brackets, not a numerical value. To determine the value one must use the concrete derivative definition, and that is where the canonical normalization \(\{Q_i,P_j\}=\delta_{ij}\) enters.

The next step is more interesting. Suppose we take the Poisson bracket of a function of the coordinate with the conjugate momentum. The lecture states the answer first and then proves it. The corresponding board frame catches the left-hand side in the act of being converted into a derivative statement.

\begin{figure}[h]
\centering
\includegraphics[width=0.56\textwidth]{lecture_09_figure_03.png}
\caption{Poisson bracket with \(P\) as a derivative}
\end{figure}

The right-hand side is still being written in the frame, but the intended statement is
\begin{equation}
\{F(q),P\}=\frac{dF}{dq}.
\end{equation}
In normalized many-variable notation the same rule is
\begin{equation}
\{F(q_i),p_j\}=\delta_{ij}\frac{\partial F}{\partial q_i},
\end{equation}
and the companion rule is
\begin{equation}
\{q,F(p)\}=\frac{dF}{dp}.
\end{equation}
The order matters. To differentiate with respect to \(q\), one brackets with \(P\) on the right. To differentiate with respect to \(p\), one puts \(q\) on the left.

A nice feature of the lecture is that this is not proved in one line. It is proved pedagogically, beginning with monomials.

First take \(F(q)=1\). Then the target statement is \(\{1,P\}=0\). Using the product rule,
\begin{equation}
\{1^2,P\}=1\{1,P\}+1\{1,P\}=2\{1,P\}.
\end{equation}
But \(1^2=1\), so the left-hand side is just \(\{1,P\}\). Hence \(\{1,P\}=2\{1,P\}\), and therefore
\begin{equation}
\{1,P\}=0.
\end{equation}

Next take \(F(q)=q\). Then
\begin{equation}
\{q,P\}=1,
\end{equation}
which is simply the defining canonical relation.

Now take \(F(q)=q^2\). The product rule gives
\begin{equation}
\{q^2,P\}=q\{q,P\}+\{q,P\}q=2q,
\end{equation}
which is exactly the derivative of \(q^2\). From here the induction to \(q^n\) is immediate, linearity extends the result to polynomials, and smooth functions are then reached by the approximation argument used in the lecture.

So Poisson bracketing with the conjugate momentum is another way of differentiating.

\section{Time Evolution as a Poisson Flow}

Only after the bracket algebra has been organized does the lecture add the dynamical postulate
\begin{equation}
\dot A=\{A,H\},
\end{equation}
for any function \(A(q,p)\) on phase space. The meaning is geometric: as the phase point moves along its Hamiltonian orbit, the value of \(A\) changes, and that rate of change is the Poisson bracket with the Hamiltonian.

This already suggests the larger pattern. A flow on phase space changes functions by Poisson bracketing with its generator. For the actual physical time flow, the generator happens to be the Hamiltonian.

To see ordinary mechanics return from the algebra, take the free-particle Hamiltonian
\begin{equation}
H=\frac{P^2}{2M}.
\end{equation}
Then
\begin{equation}
\dot P=\left\{P,\frac{P^2}{2M}\right\}=0,
\end{equation}
because a function of \(P\) has vanishing bracket with any other function of \(P\). This is momentum conservation.

For the coordinate,
\begin{align}
\dot Q
&=\left\{Q,\frac{P^2}{2M}\right\}
 =\frac{1}{2M}\{Q,P^2\} \\
&=\frac{1}{2M}\left(P\{Q,P\}+\{Q,P\}P\right)
 =\frac{1}{2M}(P+P)
 =\frac{P}{M}.
\end{align}
Thus the ordinary velocity-momentum relation appears directly from the Poisson algebra. The lecture's point is that Hamilton's equations need not be imported from outside here; they can be recovered from the bracket structure itself.

\section{Canonical Transformations from Explicit Examples}

The discussion now widens again. If the Poisson relations are the structural core of mechanics, what transformations of phase-space coordinates preserve them? These are not yet symmetries of a particular Hamiltonian. They are the allowable re-descriptions of the same system.

So let us take one canonical pair \((q,p)\) and try some transformations.

The first attempt is
\begin{equation}
Q=2q,\qquad P=2p.
\end{equation}
Then
\begin{equation}
\{Q,P\}=\{2q,2p\}=4\{q,p\}=4,
\end{equation}
so the canonical normalization is lost. Stretching both directions by the same factor does not preserve the Poisson structure.

Now try
\begin{equation}
Q=2q,\qquad P=\frac12 p.
\end{equation}
Then
\begin{equation}
\{Q,P\}=\left\{2q,\frac12 p\right\}=\{q,p\}=1.
\end{equation}
This one works. In the one-pair picture the phase-space rectangle has been stretched in one direction and squeezed in the other by the compensating factor. The area is preserved.

This is a good moment to be careful. For one \(p\) and one \(q\), area preservation is the right geometric intuition. For many degrees of freedom the full condition is more subtle, and the lecture explicitly warns against overselling the one-pair picture.

A second successful example is a rotation in phase space:
\begin{equation}
P=\cos\theta\,p+\sin\theta\,q,\qquad
Q=-\sin\theta\,p+\cos\theta\,q.
\end{equation}
By antisymmetry,
\[
\{P,P\}=0,\qquad \{Q,Q\}=0.
\]
So the only bracket worth checking is \(\{Q,P\}\). Expanding by linearity,
\begin{align}
\{Q,P\}
&=\{-\sin\theta\,p+\cos\theta\,q,\cos\theta\,p+\sin\theta\,q\} \\
&=\cos^2\theta\,\{q,p\}-\sin^2\theta\,\{p,q\} \\
&=\cos^2\theta+\sin^2\theta \\
&=1.
\end{align}
So the rotation preserves the canonical bracket relations as well.

\begin{definition}
A canonical transformation is a change of phase-space coordinates
\[
(q_i,p_i)\longmapsto (Q_i,P_i)
\]
that preserves the canonical Poisson-bracket relations
\[
\{Q_i,Q_j\}=0,\qquad \{P_i,P_j\}=0,\qquad \{Q_i,P_j\}=\delta_{ij}.
\]
\end{definition}

These transformations are not gauge redundancies. They are changes of description of the same system. One is not mapping a harmonic oscillator into some unrelated many-particle system. One is choosing new coordinates on the same phase space.

\subsection{Question \& Answer}

Question. What kinds of changes of \(p\) and \(q\) preserve mechanics, and why does area preservation enter?

Answer. For one canonical pair, the successful examples preserve phase-space area: the squeeze-stretch transformation and the phase-space rotation both do so, while the double stretch does not. But the real definition is algebraic, not pictorial. What must be preserved is the Poisson structure itself. Area preservation is the right local intuition for one pair, not the full many-dimensional definition.

\section{Infinitesimal Canonical Transformations and Generators}

The lecture now turns to a more powerful description. Instead of checking transformations one by one, we ask for a constructive description of infinitesimal canonical transformations.

Write
\begin{equation}
Q=q+\delta q,\qquad P=p+\delta p,
\end{equation}
where \(\delta q\) and \(\delta p\) are small functions of \(q\) and \(p\). Small means that quadratic and higher powers of these quantities may be neglected.

To preserve the canonical bracket to first order, we require
\begin{align}
\{Q,P\}
&=\{q+\delta q,p+\delta p\} \\
&=\{q,p\}+\{\delta q,p\}+\{q,\delta p\}+\{\delta q,\delta p\}.
\end{align}
The last term is quadratic in small quantities and is dropped. Therefore the first-order canonical condition is
\begin{equation}
\{\delta q,p\}=-\{q,\delta p\}.
\end{equation}

The lecture then gives a constructive way to satisfy this condition. Introduce a small parameter \(\epsilon\) and a phase-space function \(G(q,p)\), called the generator, and define
\begin{equation}
\delta q=\epsilon\{q,G\},\qquad \delta p=\epsilon\{p,G\}.
\end{equation}
Using the derivative form of the bracket, these become
\begin{equation}
\delta q=\epsilon\,\frac{\partial G}{\partial p},\qquad
\delta p=-\epsilon\,\frac{\partial G}{\partial q}.
\end{equation}

Now insert them into the first-order canonical condition. The left-hand side is
\[
\{\delta q,p\}
=\epsilon\left\{\frac{\partial G}{\partial p},p\right\}
=\epsilon\,\frac{\partial^2 G}{\partial q\,\partial p},
\]
while the right-hand side is
\[
-\{q,\delta p\}
=-\left\{q,-\epsilon\,\frac{\partial G}{\partial q}\right\}
=\epsilon\,\frac{\partial^2 G}{\partial p\,\partial q}.
\]
These agree identically. So a flow generated by Poisson bracketing with \(G\) is canonical.

There is an even more useful reformulation. If \(A(q,p)\) is any phase-space function, then
\begin{align}
\delta A
&=\frac{\partial A}{\partial q}\delta q+\frac{\partial A}{\partial p}\delta p \\
&=\epsilon\left(\frac{\partial A}{\partial q}\frac{\partial G}{\partial p}
-\frac{\partial A}{\partial p}\frac{\partial G}{\partial q}\right) \\
&=\epsilon\{A,G\}.
\end{align}
So the change of any observable along the flow is generated by Poisson bracketing with \(G\).

This is the right point to connect back to time evolution. If we write a small time step as \(\epsilon=\delta t\), then Hamiltonian evolution is simply the special case \(G=H\). In that sense Hamiltonian flow is itself a canonical transformation.

The lecture is cautious about how sweeping a general statement one should make here. The safe claim, and the one we keep, is this: every infinitesimal canonical transformation is generated by some \(G\), and the finite canonical transformations of practical interest are built from such infinitesimal steps.

\subsection{Question \& Answer}

Question. If \(G\) generates a flow exactly the way \(H\) does, what distinguishes a general generator from the actual Hamiltonian?

Answer. Mathematically, very little. The same formalism generates both flows. Physically, everything depends on which function is singled out by experiment as the energy of the system. A general generator may be treated as though it were a Hamiltonian for the purpose of describing a flow, but only the physics decides which one is the Hamiltonian of the actual system.

\section{Symmetry, Conservation, and the Free Particle}

We are now finally ready to define symmetry in the phase-space language. A symmetry is a canonical transformation whose flow does not change the Hamiltonian. In geometric language, the symmetry flow moves the phase point along surfaces of constant energy.

If \(G\) generates the flow, then the change of any function is \(\delta A=\epsilon\{A,G\}\). Applying this to \(A=H\), we obtain the symmetry condition
\begin{equation}
\{H,G\}=0.
\end{equation}
But by antisymmetry this is equivalent to
\begin{equation}
\{G,H\}=0.
\end{equation}
Using the Hamiltonian time-evolution rule, we immediately conclude
\begin{equation}
\dot G=\{G,H\}=0.
\end{equation}
Thus a generator whose flow leaves the Hamiltonian unchanged is itself conserved by the Hamiltonian flow.

\subsection{Question \& Answer}

Question. Why does ``\(G\) leaves \(H\) unchanged'' turn into ``\(G\) is conserved''?

Answer. Because the two statements are the same Poisson-bracket equation read in opposite orders. Saying that the \(G\)-flow leaves the energy unchanged is \(\{H,G\}=0\). By antisymmetry this is the same as \(\{G,H\}=0\). But \(\{G,H\}\) is exactly \(\dot G\). So the same vanishing bracket encodes both symmetry and conservation.

The lecture closes with a concrete example. In two dimensions take the free-particle Hamiltonian
\begin{equation}
H=\frac{p_x^2+p_y^2}{2m},
\end{equation}
and the angular momentum
\begin{equation}
G=xp_y-yp_x.
\end{equation}
Susskind briefly sets \(m=1\) during the board calculation; the factor is inessential, so we keep it explicit. We now compute:
\begin{align}
\{G,H\}
&=\frac{1}{2m}\{xp_y-yp_x,p_x^2+p_y^2\} \\
&=\frac{1}{2m}\left(
\{xp_y,p_x^2\}
+\{xp_y,p_y^2\}
-\{yp_x,p_x^2\}
-\{yp_x,p_y^2\}
\right).
\end{align}
Here the only nonzero basic brackets are
\[
\{x,p_x\}=1,\qquad \{y,p_y\}=1,
\]
with all cross-brackets vanishing. Therefore
\begin{align}
\{xp_y,p_x^2\} &= 2p_xp_y, \\
\{xp_y,p_y^2\} &= 0, \\
\{yp_x,p_x^2\} &= 0, \\
\{yp_x,p_y^2\} &= 2p_xp_y.
\end{align}
Hence
\begin{equation}
\{G,H\}=\frac{1}{2m}\left(2p_xp_y-2p_xp_y\right)=0.
\end{equation}
So the angular momentum is conserved for the free particle.

The reciprocal viewpoint is equally instructive. Suppose, not as the free particle but as a new system, we take
\begin{equation}
H=xp_y-yp_x.
\end{equation}
Then Hamilton's equations give
\begin{equation}
\dot x=\{x,H\}=-y,\qquad \dot y=\{y,H\}=x.
\end{equation}
That is circular motion. And since the Poisson bracket with the old free-particle energy vanishes, the quantity \(p_x^2+p_y^2\) is conserved in this new system. Either member of a vanishing Poisson-bracket pair may serve as the Hamiltonian; the other is then conserved.

\section{Summary}

The lecture has a very clear arc. We begin with phase space as a structured space, not merely as a list of coordinates. We then abstract Poisson brackets into a small algebra of antisymmetry, linearity, and the product rule; recover the canonical brackets and the derivative identities; add time evolution as Hamiltonian flow; and finally widen the picture to canonical transformations generated by arbitrary functions on phase space.

The payoff is the reciprocity between symmetry and conservation. A generator \(G\) defines a canonical flow. If that flow leaves the Hamiltonian unchanged, then \(\{H,G\}=0\), and the same equation says that \(G\) is conserved under time evolution. In this way the Poisson structure of phase space does not merely organize computations. It tells us what motions, transformations, and invariants classical mechanics naturally possesses.